% !TEX program = xelatex
\documentclass[11pt,a4paper]{ctexart}
\usepackage[a4paper,margin=1.7cm,headheight=15pt]{geometry}
\usepackage{fontspec}
\usepackage{booktabs,tabularx,array}
\usepackage[table]{xcolor}
\usepackage{enumitem}
\usepackage{tcolorbox}
\usepackage{fancyhdr}
\usepackage{hyperref}
\usepackage{tikz}
\usetikzlibrary{arrows.meta,positioning,shapes.geometric,fit,backgrounds}

\setmainfont{Microsoft YaHei}
\setsansfont{Microsoft YaHei}
\definecolor{navy}{HTML}{173B57}
\definecolor{blue}{HTML}{2E6F9E}
\definecolor{bluebg}{HTML}{E7F2F8}
\definecolor{green}{HTML}{2C7A5B}
\definecolor{greenbg}{HTML}{E8F4ED}
\definecolor{amber}{HTML}{B66A13}
\definecolor{amberbg}{HTML}{FFF2DA}
\definecolor{purple}{HTML}{6A4C93}
\definecolor{purplebg}{HTML}{F0EAF7}
\definecolor{red}{HTML}{A33A32}
\definecolor{redbg}{HTML}{FBE8E6}
\definecolor{graybg}{HTML}{F3F5F7}
\definecolor{graytext}{HTML}{59636D}

\hypersetup{colorlinks=true,linkcolor=navy,urlcolor=blue,pdftitle={比赛级三人滚动流水线 SOP}}
\setlength{\parindent}{0pt}
\setlength{\parskip}{4pt}
\setlist[itemize]{leftmargin=1.55em,itemsep=2pt,topsep=2pt}
\renewcommand{\arraystretch}{1.28}
\newcolumntype{Y}{>{\raggedright\arraybackslash}X}
\newcolumntype{C}[1]{>{\centering\arraybackslash}p{#1}}
\pagestyle{fancy}
\fancyhf{}
\lhead{三人建模队｜滚动流水线 SOP}
\rhead{精简执行版 V2.0}
\cfoot{\thepage}

\newtcolorbox{keybox}[1]{colback=bluebg,colframe=blue,title=#1,fonttitle=\bfseries,boxrule=.8pt,arc=2mm}
\newtcolorbox{redbox}[1]{colback=redbg,colframe=red,title=#1,fonttitle=\bfseries,boxrule=.8pt,arc=2mm}

\begin{document}

\section*{① 整体详细步骤（流程图）}

{\color{graytext}每问固定流转：①出 V0 → ②复核强化 → ③同步成文 → ①最终定版。①的 AI 或程序结果必须先经②复核。}

\begin{center}
\resizebox{.88\textwidth}{!}{%
\begin{tikzpicture}[
  node distance=5mm,
  every node/.style={font=\small,align=center},
  main/.style={rounded corners=2mm,draw=navy,fill=graybg,very thick,text width=119mm,minimum height=10mm,inner sep=3mm},
  one/.style={rounded corners=2mm,draw=blue,fill=bluebg,very thick,text width=119mm,minimum height=13mm,inner sep=3mm},
  two/.style={rounded corners=2mm,draw=green,fill=greenbg,very thick,text width=119mm,minimum height=13mm,inner sep=3mm},
  three/.style={rounded corners=2mm,draw=purple,fill=purplebg,very thick,text width=119mm,minimum height=13mm,inner sep=3mm},
  gate/.style={diamond,aspect=3.2,draw=amber,fill=amberbg,very thick,text width=45mm,inner sep=1mm},
  final/.style={rounded corners=2mm,draw=navy,fill=navy,text=white,very thick,text width=119mm,minimum height=11mm,inner sep=3mm},
  arr/.style={-{Latex[length=2.5mm]},very thick,draw=navy},
  back/.style={-{Latex[length=2.3mm]},thick,dashed,draw=red}
]
\node[main] (start) {\textbf{三人共同破题}\\明确每问输入、输出、变量、目标、约束、五问依赖、主路线与备用路线};
\node[main,below=of start] (base) {\textbf{并行建立底座}\\①搭全题架构并做 Q1｜②做数据审计、EDA、基线和验证方案｜③搭论文骨架、符号表和图表规范};
\node[one,below=of base] (v0) {\textbf{① 当前问 V0：提出并实现}\\建立模型与程序，输出初步结果、图表、结论、备用模型和风险};
\node[gate,below=of v0] (g0) {V0 是否可独立运行？};
\node[two,below=of g0] (v1) {\textbf{② 当前问 V1：独立复现并强化}\\核对数据、单位、公式和约束；补充基线、误差、收敛、敏感性或稳定性检验；筛选图表};
\node[gate,below=of v1] (g1) {结果是否可信且边界清楚？};
\node[three,below=of g1] (draft) {\textbf{③ 当前问论文初稿：证据化成文}\\按“分析—假设—建模—求解—结果—检验—结论”写作；每个结论对应图、表或程序};
\node[gate,below=of draft] (g2) {模型含义、题目回答、结论边界是否正确？};
\node[one,below=of g2] (freeze) {\textbf{① 审核并定版}\\确认模型没有写错、结果回答题目、结论没有越界；生成本问 Frozen，作为下一问唯一接口};
\node[gate,below=of freeze] (next) {是否还有下一问？};
\node[final,below=of next] (finish) {\textbf{五问完成后统一收口}\\①统一模型与正式图表｜②统一假设、检验与数字｜③统一全文、摘要、结论与排版};

\draw[arr] (start)--(base); \draw[arr] (base)--(v0); \draw[arr] (v0)--(g0);
\draw[arr] (g0)--node[right]{是}(v1); \draw[arr] (v1)--(g1); \draw[arr] (g1)--node[right]{是}(draft);
\draw[arr] (draft)--(g2); \draw[arr] (g2)--node[right]{是}(freeze); \draw[arr] (freeze)--(next);
\draw[arr] (next)--node[right]{否}(finish);
\draw[back] (g0.west) -- ++(-18mm,0) |- node[pos=.25,left]{否：退回补齐} (v0.west);
\draw[back] (g1.east) -- ++(18mm,0) |- node[pos=.25,right]{否：②否决并退回①} (v0.east);
\draw[back] (g2.west) -- ++(-26mm,0) |- node[pos=.23,left]{否：按问题性质退回②或③} (v1.west);
\draw[back] (next.east) -- ++(26mm,0) |- node[pos=.22,right]{是：进入下一问 V0} (v0.east);
\end{tikzpicture}%
}
\end{center}

\clearpage
\section*{② 每个人要做什么}

\begin{tcolorbox}[colback=bluebg,colframe=blue,title={① 总建模与技术路线负责人｜提出、实现、取舍、定版},fonttitle=\bfseries,boxrule=1pt,arc=2mm]
\begin{itemize}
  \item 共同破题后，确定五问的总路线、主模型、备用模型和前后接口。
  \item 对每一问先做 V0：定义变量、公式、约束和算法，写出可运行程序，生成初步结果与图表。
  \item 把模型、代码、结果、初步结论和风险完整交给②，不只交截图或口头结论。
  \item 根据②的复核意见修改模型；遇到数据泄漏、非法解、外推或无法复现时必须返工。
  \item 审核③的章节，只确认模型含义、题目响应和结论边界；通过后冻结该问。
  \item 五问完成后统一模型接口、正式图表、最终方案及摘要中的核心数字。
\end{itemize}
\end{tcolorbox}

\begin{tcolorbox}[colback=greenbg,colframe=green,title={② 技术复核与模型强化负责人｜复现、检验、否决、增强},fonttitle=\bfseries,boxrule=1pt,arc=2mm]
\begin{itemize}
  \item 在 Q1 V0 出来前，独立完成数据审计、变量字典、EDA、基线模型和验证方案。
  \item 收到每问 V0 后，从明确入口独立运行，核对数据、单位、公式、约束、随机性和结果一致性。
  \item 与基线比较，补充交叉验证、残差、外推、收敛、敏感性、稳定性或鲁棒性检验。
  \item 判断哪些结论有证据、哪些表述越界；决定图表保留、删除、合并或补充。
  \item 形成 V1 技术包：正式数字、模型检验、假设与理由、适用范围、可写入论文的结论。
  \item 对泄漏、非法方案、无法复现、明显外推或未收敛结果行使技术否决权。
\end{itemize}
\end{tcolorbox}

\begin{tcolorbox}[colback=purplebg,colframe=purple,title={③ 论文主笔与结构负责人｜组织、成文、追证、整合},fonttitle=\bfseries,boxrule=1pt,arc=2mm]
\begin{itemize}
  \item 比赛开始即完成问题重述、全文骨架、符号表、参考文献库和图表/公式规范。
  \item 只依据②通过复核的 V1 写作，不从 V0 或聊天记录中自行挑选数字。
  \item 每问按“问题分析—模型准备/假设—模型建立—求解—结果—检验—结论”成文。
  \item 为公式补解释，为图表补分析；明确区分观测、计算、预测、解释和情景结论。
  \item 确保每个结论都有图、表或程序结果支持，并自然衔接到下一问。
  \item 五问完成后统一全文、摘要、关键词、结论、符号、编号、引用和排版。
\end{itemize}
\end{tcolorbox}

\begin{redbox}{三人共同遵守}
文件是正式交接物；②未通过，③不写正式结论；①未冻结，下一问不得把该结果当成确定接口。
\end{redbox}

\clearpage
\section*{③ 每个人干活的节点}

\begin{center}
\small
\begin{tabularx}{\textwidth}{@{}p{2.7cm}C{1.2cm}C{1.2cm}C{1.2cm}Y@{}}
\toprule
流程节点 & ① & ② & ③ & 此节点的动作与交接\\
\midrule
共同破题 & ● & ● & ● & 共同确定问题链、变量、目标、约束、主路线、备用路线与统一口径。\\
并行建立底座 & ● & ● & ● & ①搭架构并启动 Q1；②做数据与验证底座；③搭论文骨架。三人同时工作。\\
当前问 V0 & ● & ○ & ○ & ①主责建模和程序；②可提出验证需求；③只维护结构，不写未经复核的结论。\\
V0 交接 & ● & ● & — & ①交完整技术包给②；②先检查能否独立运行，缺件直接退回。\\
当前问 V1 & ○ & ● & — & ②独立复现、检验和强化；必要时②否决，①返工，直到结论可信。\\
V1 交接 & — & ● & ● & ②把正式数字、入选图表、检验、假设和可写结论交给③。\\
当前问成文 & ○ & ○ & ● & ③写章节；②回答证据和检验问题；①只处理模型含义相关问题。\\
章节审核 & ● & ○ & ● & ①检查模型含义、题目回答、结论边界；有问题由③改文或退②复核。\\
本问定版 & ● & ○ & ○ & ①冻结本问；②确认复核结论未被改写；③确认章节与正式数字一致。\\
进入下一问 & ● & ● & ● & ①用 Frozen 接口做下一问 V0；②同步准备下一问验证；③继续滚动成文。\\
五问统一收口 & ● & ● & ● & ①统一模型和图；②统一假设、检验和数字；③统一全文、摘要、结论与排版。\\
\bottomrule
\end{tabularx}
\end{center}

\vspace{5mm}
\begin{keybox}{图例与交接关系}
\textbf{● 主责}：必须产生本节点的正式交付物；\quad
\textbf{○ 配合}：提供输入、修改或确认；\quad
\textbf{— 不介入正式产出}。\\[2mm]
每问固定交接：\textbf{① V0 技术包 → ② V1 复核包 → ③ 论文章节 → ① Frozen 定版}。
\end{keybox}

\vfill
\begin{center}
{\large\bfseries\color{navy} 完成一问，交接一问，成文一问，冻结一问。}
\end{center}

\end{document}
